\documentclass[a4paper,12pt]{article}

% --- ตั้งค่าขอบกระดาษ ---
\usepackage[top=1cm, bottom=1cm, left=1cm, right=1cm]{geometry}

% --- ฟอนต์ภาษาไทย ---
\usepackage{fontspec}
\usepackage{xunicode}
\usepackage{xltxtra}
\defaultfontfeatures{Mapping=tex-text}
\setmainfont[
    Path = ../fonts/,
    Extension = .ttf,
    BoldFont = *-Bold,
    ItalicFont = *-Italic,
    BoldItalicFont = *-BoldItalic
]{THSarabunNew}

% --- ตั้งค่าการตัดคำภาษาไทย ---
\XeTeXlinebreaklocale "th"
\XeTeXlinebreakskip = 0pt plus 1pt

% --- แพ็กเกจจัดหน้า ---
\usepackage{tikz}
\usetikzlibrary{positioning, calc}
\usepackage{graphicx}
\usepackage{svg}
\usepackage{amssymb}
\usepackage{array}
\usepackage{multirow}

% --- สร้างคำสั่งทำ Checkbox ---
\newcommand{\boxsize}{1.1em} 
\newcommand{\checkbox}{%
  \begin{tikzpicture}[baseline={([yshift=-.6ex]current bounding box.center)}]
    \draw[line width=0.8pt] (0,0) rectangle (\boxsize,\boxsize);
  \end{tikzpicture}%
}
\newcommand{\checkedbox}{%
  \begin{tikzpicture}[baseline={([yshift=-.6ex]current bounding box.center)}]
    \draw[line width=0.8pt] (0,0) rectangle (\boxsize,\boxsize);
    \draw[thick] (0.2em,0.2em) -- (\boxsize-0.2em,\boxsize-0.2em);
    \draw[thick] (0.2em,\boxsize-0.2em) -- (\boxsize-0.2em,0.2em);
  \end{tikzpicture}%
}

% --- สร้างคำสั่งขีดเส้นใต้เติมเต็มพื้นที่ ---
\newcommand{\fillrule}{%
  \leavevmode\leaders\hrule height -3pt depth 3.4pt\hfill\kern0pt%
}

\begin{document}
\pagestyle{empty}

% ================= HEADER SECTION =================
\noindent
% --- โลโก้ (จัดกึ่งกลางหน้ากระดาษแบบ Absolute) ---
\makebox[\textwidth][c]{\raisebox{-2.5cm}[0pt][0pt]{\includesvg[width=2.5cm]{../images/KMUTT_Logo.svg}}}%
\vspace{-\baselineskip}

\noindent
% --- คอลัมน์ซ้าย (Publish Box + วิชา) ---
\begin{minipage}[t]{0.54\textwidth}
    \vspace{0pt}
    % กล่องเผยแพร่
    \begin{tikzpicture}
        \node[draw, line width=0.8pt, minimum width=6.8cm, minimum height=2.5cm, inner sep=5pt] {
            \begin{minipage}{6.2cm}
                \centering
                \vspace*{-0.3em}
                \textbf{\fontsize{18}{18}\selectfont \underline{การเผยแพร่ข้อสอบ}}\\[0.8em]
                \raggedright
                \fontsize{16}{16}\selectfont
                \hspace{0.5cm}\checkbox\ \raisebox{0pt}{อนุญาต} \hspace{0.8cm} \checkedbox\ \raisebox{0pt}{ไม่อนุญาต}
            \end{minipage}
        };
    \end{tikzpicture}

    \vspace{0.3cm}

    % ชื่อวิชา (อยู่ใต้กล่องซ้าย)
    \raggedright
    {\fontsize{24}{28}\selectfont \underline{วิชา}} \\[0.6cm]
    {\fontsize{85}{85}\selectfont \textbf{MTH234}} \\[1.0cm]
    {\fontsize{50}{50}\selectfont Linear Algebra}
\end{minipage}%
\hfill
% --- คอลัมน์ขวา (ข้อมูลการสอบ) ---
\begin{minipage}[t]{0.42\textwidth}
    \raggedleft % ดันเนื้อหาไปขวาสุดของคอลัมน์นี้
    \begin{minipage}[t]{0.8\linewidth} % สร้างกล่องย่อยเพื่อจัดกลุ่มข้อความ (ขยายความกว้างขึ้นเล็กน้อยแก้ปัญหาตกบรรทัด)
        \vspace{0pt}
        \raggedright % ข้อความภายในยังคงชิดซ้าย
        \fontsize{24}{28}\selectfont
        
        \textbf{เลขที่นั่งสอบ} \rule[-2pt]{2.5cm}{0.5pt} \\[0.3em]
        
        \textbf{\underline{วันที่สอบ}} \\
        22 ตุลาคม 2568 \\[0.3em]
        
        \textbf{\underline{เวลาสอบ}} \\
        09.00 - 12.00 (3 ชั่วโมง) \\[0.3em]
        
        \textbf{\underline{กลุ่มนักศึกษา}} \\
        \checkedbox\ \raisebox{0pt}{ปกติ} \\
        \vspace{0.1em}
        \checkbox\ \raisebox{0pt}{International} \\[0.3em]
        
        \textbf{\underline{ชุดข้อสอบ (Optional)}} : 1
    \end{minipage}
\end{minipage}

\vspace{0.1cm}

% ================= STUDENT INFO BOX =================
\begin{center}
\setlength{\fboxsep}{6pt}
\setlength{\fboxrule}{0.8pt}
\fbox{%
    \begin{minipage}{\dimexpr\textwidth-2\fboxsep-2\fboxrule\relax}
        \fontsize{20}{26}\selectfont 
        \begin{tabular}{@{} p{0.62\linewidth} @{\hspace{1ex}} p{\dimexpr0.38\linewidth-1ex\relax} @{}}
            \textbf{ชื่อ-นามสกุล :} \hrulefill & \textbf{รหัสนักศึกษา :} \hrulefill \\
            \textbf{ภาควิชา/สาขาวิชา :} \hrulefill & \textbf{กลุ่ม :} \hrulefill \\
        \end{tabular}
    \end{minipage}%
}
\end{center}

\vspace{0.1cm}

% ================= INSTRUCTIONS =================
\noindent
\begin{tabular}{@{} p{2.8cm} p{\dimexpr\linewidth-2.8cm-0.5cm\relax} @{}}
    \fontsize{16}{20}\selectfont \textbf{\underline{คำชี้แจง}} & \fontsize{16}{20}\selectfont 1. ข้อสอบวิชานี้มี 6 ข้อ คะแนนเต็ม 40 คะแนน จำนวน 10 หน้า (รวมใบปะหน้า) \textbf{\underline{ให้ทำลงใน}\underline{ข้อสอบ}} \\[0.4em]
    & \fontsize{16}{20}\selectfont 2. ให้เขียนชื่อและรหัสนักศึกษาลงในข้อสอบทุกหน้า \\[0.4em]
    & \fontsize{16}{20}\selectfont 3. ให้ทำข้อสอบโดยเขียนด้วยปากกาหรือดินสอสีดำเข้ม ถ้าเนื้อที่ไม่พอให้เขียนด้านหลังของข้อนั้น ๆ \\[0.4em]
    & \fontsize{16}{20}\selectfont 4.~ข้อสอบไม่มีการแก้ไข ในกรณีที่ข้อสอบไม่ชัดเจนหรือมีข้อสงสัยให้ตัดสินใจแก้ปัญหาพร้อมทั้งอธิบายเหตุผลที่ตัดสินใจทำเช่นนั้น
\end{tabular}

\vspace{0.1cm}

% ================= CONDITIONS =================
\noindent
\begin{tabular}{@{} p{2.8cm} p{\dimexpr\linewidth-2.8cm-0.5cm\relax} @{}}
    \fontsize{16}{20}\selectfont \smash{\textbf{\underline{เงื่อนไขการสอบ}}} & \fontsize{16}{20}\selectfont 1. \vphantom{ที่จุ}\textbf{ให้นักศึกษาปฏิบัติแนวปฏิบัติการสอบอย่างเคร่งครัด} \\[0.4em]
    & \fontsize{16}{20}\selectfont 2. \vphantom{ที่จุ}\textbf{ไม่อนุญาต}ให้นำเครื่องคำนวณตามระเบียบมหาวิทยาลัยเข้าห้องสอบ \\[0.4em]
    & \fontsize{16}{20}\selectfont 3. \vphantom{ที่จุ}\textbf{ไม่อนุญาต}ให้นำตำราและเอกสารใด ๆ เข้าห้องสอบ \\[0.4em]
    & \fontsize{16}{20}\selectfont 4. \vphantom{ที่จุ}\textbf{ห้าม}ประพฤติปฏิบัติอันเป็นเหตุแห่งการทุจริตทุกรูปแบบในระหว่างการสอบโดยเด็ดขาด
\end{tabular}

\vspace{0.1cm}
\noindent
{\fontsize{16}{20}\selectfont \textbf{\underline{หมายเหตุ (Optional)}}}

\vfill

% ================= FOOTER =================
\begin{center}
    \fontsize{16}{20}\selectfont * ข้อสอบได้ผ่านการพิจารณาจากคณะกรรมการประจำภาควิชาคณิตศาสตร์แล้ว *
\end{center}

\vspace{-0.2cm}

\begin{center}
\begin{tikzpicture}
    % กล่องซ้าย
    \node[draw, line width=0.8pt, minimum width=9cm, minimum height=3.0cm] (box1) at (0,0) {};
    \node[anchor=north] at (box1.north) [yshift=-8pt] {\fontsize{16}{16}\selectfont \textbf{\underline{หัวหน้าภาควิชา}}};
    % เส้นประเซ็นชื่อ (ปรับให้เท่ากับกล่องขวา)
    \draw[dotted, thick] ([yshift=-0.3cm, xshift=-3.5cm]box1.center) -- ([yshift=-0.3cm, xshift=3.5cm]box1.center);

    % กล่องขวา
    \node[draw, line width=0.8pt, minimum width=9cm, minimum height=3.0cm, anchor=west] (box2) at ([xshift=0.5cm]box1.east) {};
    \node[anchor=north] at (box2.north) [yshift=-8pt] {\fontsize{16}{16}\selectfont \textbf{\underline{อาจารย์ผู้สอน / ผู้ประสานงานรายวิชา}}};
    % เส้นประเซ็นชื่อ
    \draw[dotted, thick] ([yshift=-0.3cm, xshift=-3.5cm]box2.center) -- ([yshift=-0.3cm, xshift=3.5cm]box2.center);
    
    % เบอร์โทรศัพท์
    % วางข้อความเริ่มที่ขอบซ้ายของเส้นประ (-3.5cm จากตรงกลาง)
    \node[anchor=west, inner sep=0pt] (phoneLabel) at ([yshift=-1.0cm, xshift=-3.5cm]box2.center) {\fontsize{14}{14}\selectfont เบอร์โทรศัพท์ : \hspace{1ex}};
    % ลากเส้นจุดจากท้ายข้อความ ไปจนถึงขอบขวาของเส้นประ (+3.5cm จากตรงกลาง)
    % ใช้ coordinate calculation (phoneLabel.base east) -- (phoneLabel.base east -| rightEdge) เพื่อให้เส้นตรงเป๊ะและอยู่ระดับเดียวกับตัวหนังสือ
    \path ([xshift=3.5cm]box2.center) coordinate (rightEdge);
    \draw[dotted, thick] (phoneLabel.base east) -- (phoneLabel.base east -| rightEdge);
\end{tikzpicture}
\end{center}

\end{document}
